% Created 2026-10-21 Wed 01:35
% Intended LaTeX compiler: pdflatex
\documentclass[11pt]{article}
\usepackage[utf8]{inputenc}
\usepackage[T1]{fontenc}
\usepackage{graphicx}
\usepackage{longtable}
\usepackage{wrapfig}
\usepackage{rotating}
\usepackage[normalem]{ulem}
\usepackage{amsmath}
\usepackage{amssymb}
\usepackage{capt-of}
\usepackage{hyperref}
\documentclass[10pt,a4paper,showtrims]
\hypersetup{colorlinks=true, urlcolor={blue}, linkcolor={blue}}
\usepackage[labelfont=bf]{caption}
\usepackage{tcolorbox}
\author{Triscuitcircuit}
\date{\today}
\title{Newt\\\medskip
\large A MacOS OpenXR/Vulkan Runtime Core}
\hypersetup{
 pdfauthor={Triscuitcircuit},
 pdftitle={Newt},
 pdfkeywords={},
 pdfsubject={},
 pdfcreator={Emacs 30.2 (Org mode 9.7.11)},
 pdflang={English}}
\begin{document}

\maketitle
\tableofcontents

\clearpage
\begin{quote}
Well she turned me into a newt.
A newt?
\ldots{} I got better.

-- Monty Pyhton and the Holy Grail
\end{quote}
\section{Newt}
\label{sec:org7f16d57}
\section{About}
\label{sec:orgd98eb26}
\texttt{Newt} is a runtime wrapper for OpenXR, MoltenVK and Monado rendering under the MIT license.
It is named after the Apple Newton, a device that was ahead of its time in many ways.

Also, a Newt is \emph{not} a lizard, and this wrapper makes disconnecting from a certain walled garden just a bit easier.

Newt is split into two systems:
\begin{itemize}
\item \texttt{Newt} ((\texttt{libnewt\_core.a}): A runtime library for communicating with Monando, OpenXR and MoltenVK
\item \texttt{newtd}: A session-manager daemon for transfering to game engines
\begin{itemize}
\item (A few Godot examples are included in this repository, which include usage of the \texttt{Newt} bridge addon)
\end{itemize}
\end{itemize}
\subsection{Capabilities}
\label{sec:org99c01cb}
\begin{itemize}
\item OpenXR session lifecycle, controllers, and haptic output via the \texttt{Newt::Session} C++ API.
\item GPU (IOSurface) texture import to Vulkan, with media capture built on top of it.
\item The daemon, \texttt{Newtd}, handles one-owner headset supervision.
\item An optional generic libusb module (\texttt{NEWT\_WITH\_USB=ON}).
\item Initiate headset pairing on macOS.
\end{itemize}
\section{Installation}
\label{sec:orgb948b55}
Building Newt and configuring its headset runtime are separate steps. The OpenXR SDK supplies the \emph{loader}. The headset backend supplies a \emph{runtime}. Installing the loader does not create Newt's runtime configuration.
\subsection{Requirements (from source build)}
\label{sec:org7925b9e}
\begin{itemize}
\item MacOS, with Xcode Command line tools (\texttt{xcode-select -{}-{}install})
\item CMAKE 3.20+ ( \href{https://cmake.org/}{CMAKE.org} or by \href{https://formulae.brew.sh/formula/cmake}{Homebrew})
\item The \href{https://vulkan.lunarg.com/sdk/home}{Vulkan SDK} for MacOS, including \texttt{glslangValidator}
\item An installed \href{https://github.com/KhronosGroup/OpenXR-SDK}{Khronos OpenXR SDK} and loader, with its CMake package
\item \emph{Optional:} libusb + pkg-config, detected automatically ( \texttt{brew install libusb pkg-config} )
\item A compatible external OpenXR runtime (Monando)
\item Python (for automatic tests)
\end{itemize}
\begin{quote}
{[}!NOTE]
Do not set \texttt{VULKAN\_SDK} to \texttt{libMoltenVK.dylib}. Newt needs the SDK's \texttt{macOS} directory (headers, libraries, shader tools). Also, Newt's setup expects a runtime JSON manifest instead of just \texttt{libopenxr\_loader.dylib}.
\end{quote}
\subsection{Building and Installation}
\label{sec:org9639605}
The CMAKE build process is designed to find the SDKs. Specifically, it finds Vulkan with \texttt{/VulkanSDK} and OpenXR SDK under your \texttt{brew} install path. (\texttt{/opt/homebrew} or \texttt{/usr/local}).
\begin{verbatim}
cmake -S . -B build -DCMAKE_BUILD_TYPE=Release -DCMAKE_INSTALL_PREFIX=/your/install/prefix -DBUILD_TESTING=ON
cmake --build build
ctest --test-dir build --output-on-failure
cmake --install build
\end{verbatim}
This builds a static library of Newt (\texttt{libnewt\_core.a}) and the native \texttt{newtd} session-management Command Line Interface
\subsubsection{Optional Commands}
\label{sec:orgee35afe}
\begin{itemize}
\item \texttt{NEWT\_WITH\_USB=AUTO} for building with Libusb.
\item \texttt{NEWT\_BUILD\_DAEMON=OFF} to skip building \texttt{newtd} entirely.
\end{itemize}
\subsubsection{Linking with CMAKE}
\label{sec:org0ce0959}
If you want to link against Newt with CMAKE, you can use  \texttt{find\_package(Newt CONFIG REQUIRED)} with \texttt{Newt::core}. .

If OpenXR isn't under one of those curated prefixes, you can direct the build with: \texttt{-DOpenXR\_DIR=/path/to/folder/containing/OpenXRConfig.cmake} for one folder, or \texttt{-DNEWT\_SDK\_SEARCH\_PATH=/some/prefix} to search a whole path. You can also use the included Python code to find likely install locations.
\begin{verbatim}
python3 tools/find-libraries.py --root /path/to/sdk-sources
\end{verbatim}
\subsection{Configuring the runtime (one time)}
\label{sec:org7f6f53b}
Set \texttt{XR\_RUNTIME\_JSON} to the runtime manifest and configure the Vulkan loader (if the runtime requires it).
\begin{verbatim}
newt configure --runtime /path/to/runtime.json \
               --icd /path/to/MoltenVK_icd.json
# or, for a host provided dashboard executable (like PippinVR):
newt configure --dashboard /path/to/host-dashboard \
               --runtime /path/to/runtime.json \
               --icd /path/to/MoltenVK_icd.json
\end{verbatim}
This can also be done when installing with CMAKE by passing  \texttt{-DNEWT\_CONFIGURE\_RUNTIME=/path/to/runtime.json} and the \texttt{cmake -{}-{}install} runs the equivalent \texttt{newt configure}.
Configuration is stored in \texttt{\textasciitilde{}\textasciitilde{}/Library/Application Support/Newt/config.json}. From a source checkout, use the built daemon (\texttt{build/newtd}, installed as \texttt{newt}) in place of the installed \texttt{newt} command.

In Godot, the same paths can be picked with \textbf{Project -> Tools ->  Configure Newt…}.
\section{Embedding Newt}
\label{sec:org2dc1f56}
\subsection{The C++ API}
\label{sec:org9beed48}
As mentioned you can link \texttt{Newt::core}, and include \texttt{<Newt/Session.h>} in your code. \texttt{Newt::Session::Run(delegate)} owns the initialization process input polling, frame submission, and cleanup, synchronously, on the calling thread.

\texttt{RequestStop()} may be called from another thread. The delegate hooks are:
\begin{itemize}
\item \texttt{OnSystem} (host policy checks)
\item \texttt{OnReady} (allocate host resources)
\item \texttt{OnInput}, \texttt{OnPrepareFrame} (upload content)
\item \texttt{OnDrawEye} (draw inside an active render pass)
\item \texttt{OnStopping} (release resources before the device is destroyed).
\end{itemize}

Logging can be routed to the host through \texttt{include/NewtLog.h}. However, if no logging handler is installed, the output will go to stdout.
\section{Design}
\label{sec:orgfa898cb}
Newt is split into a product-neutral runtime core and into different layers off the core.

\begin{center}
\begin{tabular}{ll}
Component & Responsibility\\
\hline
Core (\texttt{newt\_core}, \texttt{Newt::core}) & OpenXR session lifecycle, tracking/controller input, haptic output, Vulkan resources, swapchains, frame submission, GPU (IOSurface) texture import.\\
Session manager (\texttt{daemon/newtd.mm}, binary \texttt{newtd}) & One-owner supervision of a game and a host's dashboard,Game engine handoff, developer and standalone bundles.\\
External runtime (e.g. Monado + MoltenVK) & Headset USB transport and graphics translation to Metal.\\
\end{tabular}
\end{center}


OpenXR swapchain images are borrowed and never destroyed by Newt.
\subsection{Known Issues}
\label{sec:orgde0f8b8}
\begin{itemize}
\item Every dashboard to  game handoff (and vice versa) requires different USB pairing headset requests.
\item Headset USB audio is not implemented.
\item A Godot game launched through the manager loads two independent MoltenVK copies into one proces.
\end{itemize}
\section{Testing}
\label{sec:org7161a33}
Build with \texttt{BUILD\_TESTING=ON} (the installer does), then:
\begin{verbatim}
ctest --test-dir build --output-on-failure
python3 tests/session_manager_test.py
\end{verbatim}
The suites cover session lifecycle, the session-manager protocol against fake children (including the no-dashboard-configured path), optional-USB error cleanup, and a core-boundary check (fresh install, consumed by a plain C++ program with no product UI anywhere in the tree). None of them start VR, open a camera, or claim a USB device. The manager suite needs permission to create local Unix sockets.
\section{Contributing}
\label{sec:org194cf40}
Contributions are welcome, and keep Newt free of vendored binaries — see \texttt{DEPENDENCIES.md}.

This project will take contributions in the form of pull-requests.
\section{License}
\label{sec:org0c8fbd7}
MIT License. See \texttt{LICENSE}.
\subsection{Libraries used}
\label{sec:org0043671}
\begin{itemize}
\item OpenXR SDK/loader (Apache-2.0/MIT)
\item Vulkan/MoltenVK (Apache-2.0)
\item (optional and dynamically linked) libusb (LGPL-2.1-or-later)
\end{itemize}
\end{document}
